\chapter{Integrating {\meq} into a modelling framework}
\label{ch:embedding}

{\meq} is a library first and a program second. Everything the program does is
done by calling the library, and the library does a great deal the program never
touches: the flux surfaces, the flux-surface averages and the smooth
representation of the whole family of surfaces are reachable only this way, and
none of them appears in any output file.

This chapter is for somebody writing a code that calls {\meq} --- a transport
solver that needs geometry, an optimiser that needs many nearby equilibria, a
design loop, a coupled simulation. It is organised as the sequence such a code
follows: set up a solver, give it the physics, solve it, warm-start the next one
from it, and read the geometry off the answer.

\section{Getting at the library}

\subsection{Include the headers you use}

There is an umbrella header, and \emph{it does not cover the library}. It brings
in the configuration parser, the solver and the source and profile interfaces,
and it omits the error estimator, the boundary shapes, the field and sampling
helpers, the output writers, the warm-start machinery and every one of the five
headers that make up the flux-surface stack. A consumer who includes it and
stops gets less than half of {\meq} with nothing to say so.

Include what you use, by name. The headers are the primary documentation of this
library and are unusually complete; where this chapter and a header comment
disagree, the header is right and this chapter is stale.

\subsection{Everything borrows}

\begin{quote}
	\textbf{Everything handed to a solver is borrowed and must outlive it}
	--- the mesh, the source, every profile and coefficient, and any transfer
	path. Nothing here takes ownership, and
	the flux-surface objects go further: the critical-point finder and the
	contour tracer hold the solver's own fields by reference, and the solver must
	have been solved before either is used.
\end{quote}

There is one exception that matters, and it goes the other way. A source given
in normalised flux is \emph{mutated} by the solver, once per residual
evaluation, because the flux on axis is an unknown the solver is driving. Such a
source must not be shared with anything that reads it in the meantime.

\section{The solver}

\subsection{Construction, and the one hard ordering rule}

A solver is constructed on a mesh, a polynomial degree and a stabilisation. The
three finite element spaces are built there and then --- which is why a refined
mesh needs a \emph{new} solver rather than an update, and why an adaptive loop
built on this library constructs one per cycle.

The forms, by contrast, are built on the first solve rather than in the
constructor, because the linear and semi-linear paths need structurally
different ones and hybridization takes what it finds at the moment it is
enabled. That gives the single ordering rule a caller must respect:

\begin{quote}
	\textbf{The source and the boundary extension must be set before anything
	assembles.} Both refuse a late call --- they throw rather than being
	silently ignored --- and setting a second source of any kind is refused
	outright, because one solver holds one source.
\end{quote}

Every other setter is safe at any time, and they differ only in how much they
throw away. Changing the globalisation, the nonlinear ordering, the assembly
mode or the element-local solver discards the assembled forms; setting or
clearing an initial guess discards only the reduction to the trace system;
setting the boundary data, the Newton controls or the Picard parameters discards
nothing. Choosing the trace solver discards nothing either, deliberately and
unlike every other setter in that group, because the trace solver is chosen when
the reduced system is solved rather than when the forms are assembled.

That distinction is not pedantry. It was once got wrong in exactly the place it
mattered: changing the globalisation used to discard the reduction and not the
forms, while the decision about whether the potential block goes on the linear
or the nonlinear form reads the globalisation --- so a solver switched between a
Picard path and a Newton one quietly reused the other path's blocks. It was
latent for as long as every caller built a fresh solver per globalisation, and it
surfaced the moment one globalisation switched twice inside a single solve.

\subsection{Solving, and what comes back}

\texttt{solve()} prepares the system, dispatches to whichever nonlinear path is
configured and fills a residual history. Afterwards the potential, the flux and
the trace are available as grid functions; a further call post-processes, giving
the enriched potential $\psi^{\star}$ in the space one degree up, an enriched
flux, and the total flux in a divergence-conforming space. Diagnostics ---
iteration counts of every kind, the number of trace unknowns, the factorisation
counts --- are accessors, and error norms against a supplied coefficient are
there for the benefit of convergence studies.

\texttt{prepare()} is public on its own, and for one reason: it stops after
assembling and reducing, leaving the reduced operator, its right-hand side and
its gradient exposed. That is what lets a caller difference the assembled
residual against the assembled Jacobian, which is the one check that separates a
wrong Jacobian from a wrong discretisation. If you write a source, do this.

\subsection{Failure}

\begin{quote}
	\textbf{A solver is not reusable after a failed solve.} With exceptions
	enabled in the underlying library, a failure throws from the middle of it and
	leaves its objects as the throw found them. Build a fresh solver for the
	retry; the program's own fallback ladder does exactly that.
\end{quote}

Nor is a solver reusable across meshes, for the reason given above.

\subsection{Two conventions that will bite}

\begin{description}

\item[The assembled flux block holds $-q$.] The underlying mixed form is built
	for a constitutive law of the opposite sign, and the integrators that make
	the hybridization consistent have that sign baked in. The negation is undone
	once, into the grid function the flux accessor returns. So the accessor gives
	{\meq}'s $+q$ and is what the field helpers, the sampler and the contour
	tracer want; the raw block is what the transferred-datum reconstruction
	wants, and feeding it the accessor returns the datum with its sign reversed
	rather than a small bias. The total flux is deliberately left in the
	library's convention, because the constraint equation it is projected
	through is written that way.

	The failure modes of a sign error here are worth knowing because they are so
	unequal. The contour tracer survives one visibly --- it traces the same curve
	backwards. The flux-surface averages \emph{cannot see it at all}, since only
	$|q|$ enters, which is exactly why those numbers must never be quoted as
	evidence about a sign. And the critical-point finder silently swaps every
	maximum for a minimum while its own topological audit still passes.

\item[The trace is not the transferred datum.] On the curved
	path the trace degrees of freedom on $\Gamma_h$ remain in the essential list
	at zero, and the datum actually imposed is never stored anywhere. A
	dedicated accessor rebuilds it on request, which is what the error estimator
	needs and what anybody inspecting the boundary condition should use.

\end{description}

\section{The physics input}

\subsection{The source interface is two functions}

\begin{verbatim}
    class Source
    {
    public:
      virtual double f( double r, double z,
                        double psi ) const = 0;

      virtual double dFdPsi( double r, double z,
                             double psi ) const = 0;
    };
\end{verbatim}

That is the whole of it. Plain doubles, no finite element types --- which is what
makes a source unit-testable without the library --- and no $1/r$, because the
weight belongs to the weak form rather than to the physics. The vertical
coordinate is in the signature because a manufactured solution may use it; a
physical magnetohydrodynamic source does not.

\begin{quote}
	\textbf{\texttt{dFdPsi} must be the exact derivative of \texttt{f}.} An error
	there does not change the converged answer. It wrecks --- or, worse, merely
	slows --- the convergence to it, and \chref{ch:intro} explains why no
	convergence table can see that. Difference it against \texttt{f} in a unit
	test, and watch the observed order the solver prints.
\end{quote}

\subsection{Profiles have three derivative levels, and the third is not optional}

A profile supplies its value, its first derivative and its second. The third is
a pure virtual rather than something defaulting to zero, so that every
implementation has to answer rather than one silently returning nothing.

It is there because a rotating plasma's source is \emph{already} a
$\psi$-derivative of quantities built from flux functions, so its Jacobian spends
a second derivative of every input, and no reparametrisation avoids that: the
pressure of a rotating plasma is not a flux function, so there is no product to
pre-store. The static magnetohydrodynamic source dodges it by storing the
\emph{products} $p'$ and $gg'$ rather than $p$ and $g$ --- precisely so that no
chain rule exists to get wrong --- and a normalised version of the same source
differs from it by exactly one factor of $1/\psiax$, which is the chain rule and
is the whole difference.

\begin{quote}
	\textbf{A profile must not throw outside its range.} A Newton iterate
	routinely overshoots the physical interval, and an exception out of a
	quadrature loop would kill a solve that was about to converge. Clamp
	instead. The tabulated profile clamps to its end knots with both derivatives
	set to zero.
\end{quote}

One caveat on the tabulated profile, recorded because it is unmeasured rather
than because it is known to matter: a Hermite cubic is only once continuously
differentiable, so its second derivative is piecewise linear and jumps at every
interior knot. That is a set of measure zero and does not move a converged
answer, but a Newton step taken exactly at a knot sees a one-sided Jacobian, and
nothing has measured what that costs.

\subsection{Normalised flux makes the axis an unknown}
\label{sec:normalised}

If the profiles are functions of $\Psi = \psi/\psiax$ rather than of $\psi$, then
$\psiax$ is a functional of the solution and the source must be a
\texttt{NormalisedSource}: a \texttt{Source} that can additionally be told its
normalisation. The solver calls that setter once per residual evaluation, so it
has to be cheap and must not allocate.

The reason this is a different class rather than a flag is that three cheaper
alternatives were tried and each failed in an instructive way.

\emph{Fixing $\psiax$ does not approximate the problem, it replaces it.} With the
normalisation frozen at a value that is not the answer, a peaked profile
contributes essentially nothing: pressure amplitudes differing by more than two
orders of magnitude gave solutions agreeing to every digit printed, because the
normalised flux never got far from zero.

\emph{The self-consistent problem is not the parametrised one.} At a fixed
$\psiax$ this equation has both a small positive solution and a large one. Only
the large one can satisfy the constraint that closes the normalisation --- but
the small one is well within the iteration's reach, and Newton from the Dirichlet
datum walks straight onto it. This path therefore needs an initial guess of about
the right height, and that is part of the problem statement rather than a
tuning nicety.

\emph{Closing the loop outside the solver does not work, and the obstruction is a
pole.} Iterating the normalisation against the peak flux, with relaxation,
converges --- to a degenerate fixed point where both shrink together and the
pressure gradient runs away. Mapped afterwards, the cause is plain: the outer map
has a pole a fraction of a per cent away from its own fixed point. Relaxing
harder is not a fix for a pole.

What works is putting $\psiax$ inside the residual, as a bordered Newton: the
hybridized trace residual and the constraint that the normalisation equals the
peak flux, solved together with the Jacobian bordered by one dense column and one
sparse row. Those two borders are the non-local terms that a normalised profile
is known to introduce, and the reason they must live in a border rather than in
the matrix is structural --- hybridization eliminates flux and potential
\emph{element by element}, and a term coupling every element to the one element
holding the axis is precisely what that elimination cannot represent. It costs
one extra backsolve per Newton step and no second matrix.

\section{Warm starting between nearby equilibria}

An optimiser or a coupled solve moves between equilibria that are close to one
another. Doing that well is what makes {\meq} affordable in such a loop, and
there are three distinct routes with different costs.

\begin{description}

\item[Exact restart.] Same mesh, same degree. Read the stored grid function and
	hand it to the solver as an initial guess; nothing else is needed and the
	restart is bitwise.

\item[Full-order transfer.] Different mesh or different degree, same code. This
	is what \texttt{meq::FieldTransfer} is for.

\item[Interchange.] Another code entirely, through the gridded output.
	\textbf{This one throws most of the previous solve away}: bilinear
	interpolation on a structured grid is second order however fine the grid is,
	so restarting a cubic solve through it discards nearly everything the
	previous solve computed. It still converges, because a guess is only a guess
	--- but `warm' is then doing less work than it looks.

\end{description}

\subsection{The transfer object}

\texttt{FieldTransfer} is constructed on a source mesh and transfers one grid
function at a time onto a target, taking a fallback coefficient to evaluate at
any node the source mesh does not cover, and returning the number of misses.

It is batched by design and is deliberately not a coefficient: the underlying
point search amortises a hash and an exchange over the whole query, so asking for
one point at a time pays all of that per point. Setup is once per source mesh and
is the expensive part --- keep one object and move the potential and both flux
components through it.

Two diagnostics are worth reading rather than ignoring. The \emph{miss count},
because a restart that silently found no data over most of the domain has quietly
become a cold start, and the iteration count will not obviously say so. And the
\emph{worst distance}, which is subtler: large values mean the two meshes
disagree about where the domain is, which a miss count of zero will not tell you.
On the curved path there genuinely are uncovered nodes, because the computational
domain \emph{grows} as it refines.

\subsection{The trap that makes a good guess fail}

\begin{quote}
	\textbf{A convergence target scaled to the initial residual makes a good
	initial guess fail, and the better the guess the more certain the failure.}
\end{quote}

The usual Newton stopping rule is a tolerance relative to the residual at the
iterate the solver was handed. Warm-start from a converged answer and that
residual is already small, so the target shrinks with it --- and past a point it
falls below the round-off floor, where nothing can meet it. Measured on a
benchmark restarted from its own exact answer, the residual reached the floor in
two iterations and was then reported as a failure at the thirtieth.

{\meq} therefore takes its reference from the \emph{cold} iterate --- the
Dirichlet datum alone, where the solve would have started with no guess --- and
sets a pure absolute target from it. The relative tolerance keeps exactly the
meaning it always had, a cold solve is bit-identical, and only a warm one
changes: from failing to converging in a step or two. It costs one extra residual
evaluation, and only when a guess was set.

This is the failure mode that matters most for how {\meq} will be used, which is
why it is described at length rather than listed.

\subsection{What a guess does and does not seed}

A guess arrives as a coefficient, which cannot be differentiated, so the
potential and the trace are seeded and \emph{the flux block is left at zero}. On
the default nonlinear path the flux is an unknown, so the seeded state is
inconsistent in exactly the row that couples them, and the initial residual goes
\emph{up} rather than down. The guess still works --- that row is linear in the
flux, so one Newton step recovers the flux belonging to the guessed potential ---
and the measured benefit is fewer iterations and an unmoved answer, which is what
a warm start must deliver. But do not read the initial residual as evidence that
a guess took.

\section{Flux surfaces}

Almost everything downstream of an equilibrium wants the \emph{inverse} of what
the solver computes: not $\psi$ at a position, but the surfaces themselves,
indexed by flux label. {\meq} extracts them in three layers, each usable on its
own.

\subsection{Why this is cheap here}

Every step is more accurate than it would otherwise be for one reason, and it is
\chref{ch:intro}'s first commitment: $\gradbar\psi = rq$ with $q$ solved for.

A critical point is where $q$ vanishes --- a residual converging at the
potential's own order, rather than a differentiated potential converging an order
down. The tangent to a contour is a rotation of $q$, available at every point the
tracer has visited at no cost, since it evaluated $q$ for the predictor anyway.
The weight in every flux-surface average is $1/(r|q|)$, pointwise. And the fit of
\S\ref{sec:disc} needs the gradient of the normalised flux, which is the same
field once again.

The contrast is sharp. A continuous Galerkin code at lowest order confines the
magnetic axis and the X-point to mesh vertices, because that is where a discrete
gradient can be said to vanish, and CEDRES++ records this as an open problem
\cite{Heumann2015}. {\meq} resolves both inside an element, by root finding in
the element's own polynomial.

\subsection{Critical points}

The finder locates the magnetic axis and any saddle by Newton inside one element
at a time. Classification is free and needs no second derivative of $\psi$:
differentiating $rq = \gradbar\psi$ gives a Hessian whose first term vanishes
identically wherever $q$ does, so the Hessian there is $r$ times the Jacobian of
$q$ --- and $r > 0$ throughout an axisymmetric domain, so the signs that decide
maximum, minimum or saddle read straight off that Jacobian.

That Jacobian is taken by central differences, and it looks like a shortcut and
is not one. The located root is where $q_h$ vanishes, which is a property of
$q_h$ alone; the Jacobian steers the iteration and appears nowhere in its fixed
point. A wrong Jacobian costs Newton steps and buys no error --- the same
observation \chref{ch:intro} makes about the Grad--Shafranov Newton, seen from
the other side.

Three warnings.

\begin{quote}
	\textbf{{\meq}'s $\psi$ is not sign-normalised across sources, so the axis is
	not always a maximum.} With a single-signed negative source --- which is what
	the standard analytic benchmarks have --- $\psi$ is a subsolution, its
	maximum is on the boundary, and the axis is an interior \emph{minimum}. The
	finder therefore seeds from both nodal extremes and refuses rather than
	guessing if both yield a genuine interior extremum.
\end{quote}

\begin{quote}
	\textbf{This is not the solver's own axis flux, and the two must not be
	reconciled.} The solver's $\psiax$ is the largest \emph{nodal} value of the
	discrete potential, chosen deliberately so that the bordered Newton of
	\S\ref{sec:normalised} has a constraint it can differentiate. The finder's is \emph{the place
	where the flux vanishes}. They differ in position at first order in the mesh
	size and in value at second, and \emph{both orders are independent of the
	polynomial degree} --- so on a refined high-order mesh the two readings
	separate rather than converge. Neither is a defective version of the other.
	Use the solver's for the normalisation and the finder's for the geometry.
\end{quote}

\begin{quote}
	\textbf{The topological audit is a certification and never an exclusion.} The
	sum of the indices of the located points can be compared against the domain's
	Euler characteristic, and agreement is real evidence; disagreement is real
	evidence too. But a degree of zero does not imply that nothing was missed ---
	spurious critical points appear in \emph{pairs} whose indices cancel, and the
	complementary check against a persistence threshold is not implemented. Nor
	is the sweep exhaustive: it is seeded Newton, so do not read its result as
	`these are all of them'.
\end{quote}

\subsection{The tracer}

A surface is traced by prediction along the tangent and correction back onto the
level set --- the standard predictor--corrector walk \cite{AllgowerGeorg2003},
with three things chosen for this problem.

\emph{The predictor is first order and that is not a defect.} The corrector is a
root find, so every accepted point sits on the discrete level set to tolerance
whatever the predictor did; there is nothing for a higher-order predictor to buy
in accuracy. What it would buy is step length. And what the corrector buys is
that error off the contour does not accumulate along it.

\emph{The step controller equidistributes turning, not effort.} Continuation
software keys its step length on how hard the corrector worked, which is right
when the object is to get to the end of a path. Here the object is a well-shaped
\emph{curve}, so the step is chosen to equidistribute interpolation error, which
means turning through a fixed angle per step.

\emph{The interpolant between points is a cubic Hermite whose tangents come from
$q$}, which is fourth order in the step where a chord is second. The chordal
version is kept as a live control.

The resulting error has three independent parts and {\meq} reports them as three
numbers, because conflating them is how this problem gets done badly: the
\emph{field} error, from the level set of $\psi_h$ differing from that of $\psi$;
the \emph{point location} error, which is the corrector's own and does not
accumulate; and the \emph{representation} error of the interpolant between
points, which is decoupled from the mesh and the degree entirely.

\begin{quote}
	\textbf{Root the post-processed potential, which is the default.} Pairing the
	tracer with $\psi^{\star}$ rather than $\psi_h$ preserves the interpolant's
	fourth order, where the raw pairing loses it. The reason is not the obvious
	one: the local post-processing is driven by the \emph{reconstructed total
	flux} \cite{Stenberg1991}, so the gradient of $\psi^{\star}$ and $rq^{\star}$
	are not the same object --- they agree one order better than $\gradbar\psi_h$
	and $rq_h$ do, and it is that agreement the cubic Hermite needs, because the
	tangent it is given is not the tangent of the curve it is interpolating. The
	tracer refuses a solver that has not been post-processed, because the
	accessor would otherwise hand it an empty field without complaining.
\end{quote}

Some further behaviour a consumer should know about. A trace that leaves the
mesh returns a status saying so, and \emph{returns that status rather than a
shorter curve} --- the earlier code in this tree printed a message and returned a
partial arc labelled as a closed contour, and a flux-surface average over such an
arc is wrong by an amount nothing reports. A corrector that stalls is being
honest rather than failing: the level set of a discontinuous field is a union of
per-element arcs offset by the jump at each face, so a point landing within that
jump of a face may be on neither, and no point can be closer to `the' level set
than the field's own ambiguity there. The extraction is \emph{not} exhaustive ---
a disjoint island at the same level is not found, not reported and not looked
for. And it is \emph{not} X-point aware: the level set through a saddle is not a
curve and has no tangent there.

On the curved path the outermost closed surface \emph{is} the plasma boundary,
which lies in the band outside the mesh --- so the band is the interesting part
rather than a corner case, and it is where the averages are most wanted and least
forgiving. The tracer will continue into it, by either of two constructions: a
Taylor step in the flux, which is second order whatever the degree; or the
transfer technique's own path lifting, which is in principle better than the
potential's own order and so does not limit the contour at all. Both are kept, the
weaker one as the control --- if the control ever converges as fast as the answer,
the comparison is empty and the measurement is worthless. Refusing to leave the
mesh at all is the default.

\subsection{The angle parametrisation, and a trap made into a measurement}
\label{sec:angleparam}

The third layer re-expresses a traced contour as a radius against poloidal angle
about the axis. That gives a quadrature rule, and the rule needs the derivative
of the radius with respect to the angle.

\begin{quote}
	\textbf{That derivative is available pointwise from $q$, and differencing it
	instead costs an order.} Writing $x(\theta) = a + \rho(\theta)u(\theta)$ with
	$u$ the unit ray and $t$ the unit tangent, the level-set condition gives
	$\rho' = \rho\,(u\cdot t)/(u'\cdot t)$ directly, with nothing differenced and
	nothing differentiated.
\end{quote}

The class computes the wrong answer as well as the right one, deliberately: it
reports the length from the pointwise metric, the length from a second-order
differenced metric, and the naive chordal length. The second of those is the trap
--- the same spectrally accurate quadrature rule fed a second-order Jacobian,
with nothing in its own output to say so --- and it is a first-class member
rather than test scaffolding for exactly that reason.

Star-shapedness about the axis is a \emph{hypothesis}, not an assumption. It is
reported --- as the minimum transversality of ray against tangent, which is the
denominator above --- and it is checked a second, independent way, against the
traced curve's own polar angle, so that the check does not depend on the fit
having succeeded.

\section{Flux-surface averages}

This is what a transport code actually reads.

\subsection{The convention is part of the definition}
\begin{equation}
	V'(\psi) = \oint \frac{2\pi R\,\mathrm{d}l}{|\gradbar\psi|},
	\qquad
	\langle X \rangle_\psi = \frac{1}{V'}
		\oint \frac{2\pi R\, X\,\mathrm{d}l}{|\gradbar\psi|}.
	\label{eq:fsa}
\end{equation}

\begin{quote}
	\textbf{The $2\pi R$ is in.} A per-unit-length $V'$ differs by exactly that
	factor, so anything reading a $V'$ out of {\meq} is reading the volume
	derivative $\mathrm{d}V/\mathrm{d}\psi$. The identity of
	\S\ref{sec:identity}, which is this facility's acceptance criterion, is
	stated for this convention and not for the other.
\end{quote}

A second convention is worth stating because it prevents a very expensive
confusion. \emph{The safety factor is called \texttt{safetyFactor} and is never
called \texttt{q}.} In {\meq}, $q$ is the flux, a solved unknown. The safety
factor is also universally written $q$. A reader who meets \texttt{q} anywhere in
this library is entitled to assume the flux, and one silent conflation would be
very hard to find afterwards.

\subsection{One facility, not a list of quantities}

The primitive is: given a surface, take a callable integrand in position, flux
label and flux, and return its average and $V'$. Everything named --- the inverse
square radius, the squared gradient over squared radius, the arc length, the
safety factor --- is a one-line wrapper on it.

That shape is a response to the consumer rather than a preference. The list of
geometry slots a coupled transport code wants is negotiated with the physics case
rather than fixed, so a list that is still moving must not be baked in as a list
of functions. Adding or removing one is a line, the convergence machinery is
written once, and --- most importantly --- the derivative of
\S\ref{sec:jacobian} is taken of the facility rather than of eight separate
expressions.

\subsection{The order is the flux's, not the potential's}

\begin{quote}
	This was the surprise of the stage that built it, and it is worth stating
	before anybody plans around the opposite.
\end{quote}

The tracer roots $\psi^{\star}$, which converges a full order faster than
$\psi_h$, so the natural expectation is that the averages inherit the better
order. \emph{They do not.} The level set improves and the \emph{weight} does not:
the weight divides by $|\gradbar\psi| = r|q|$, and the enriched flux converges at
the same order as the raw one. The local post-processing buys its extra order in
the \emph{potential}, and there is no better flux to be had; an average built on
both inherits the worse.

So every quantity here converges at the flux's order with either pairing, and
what $\psi^{\star}$ buys is a \emph{constant} --- a factor of a few. A constant
is not an order, and the two must not be confused. This is the same shape as the
band continuation of the magnetic field in \chref{ch:using}: a quantity limited
by the one factor with no solved variable behind it.

\subsection{The metric trap does not go away because a quantity is an average}

An average is a \emph{ratio} of two integrals over the same metric, so a bad
metric looks as though it ought to cancel. It cancels a constant and nothing in
the order: with the metric differenced rather than taken pointwise, the averages
revert to second order in the node count exactly as an arc length does, and the
two columns separate by orders once the node count is large. The ratio buys a
smaller constant and no rate.

The differenced control is therefore a first-class member of the result. It is
filled only by the builder whose parametrisation the trap lives in, and asking
for it when it is absent \emph{throws} rather than quietly returning the
pointwise answer --- which would make the control agree with the answer for the
worst possible reason.

\subsection{Two extractions, and why both are kept}

There are two builders, and they share the field and the level and nothing else.
One lays out nodes equispaced in poloidal angle, each found by a
one-dimensional ray Newton, with the pointwise metric above and the periodic
trapezoidal rule \cite{TrefethenWeideman2014} --- which on a smooth periodic
integrand converges faster than any algebraic order. The other puts
Gauss--Legendre nodes on the traced curve's own cubic segments, wherever the step
controller happened to place them, with the interpolant's own metric, and assumes
no star-shapedness at all.

\emph{They agree or one of them is wrong}, and agreement is worth more than
either being plausible on its own. A missing factor of $2\pi R$, a metric taken
about the wrong point or a gradient on the wrong side of the division would each
break it, and no single-route table could see any of them.

Two honest qualifications. The two agree only to about their own error and cannot
do better, since the level set of a discontinuous field is a union of arcs and
two routes placing nodes differently sample different arcs --- so the property to
assert is the \emph{rate} at which the gap closes, not that the gap is small. And
a third leg exists in principle --- a quadrature rule built on the level set with
no curve extracted at all --- and is \emph{not} built. Read the agreement as two
legs of three.

There is also a free check that is \emph{not} a third extraction, and it is the
first appearance in this chapter of the gauge argument that dominates
\S\ref{sec:disc}. The angle builder lays its nodes out about a ray \emph{origin},
and the integral cannot depend on it: a reparametrisation of the same curve
integrates to the same number. Fitting the same contour about a deliberately
displaced origin and getting the same $V'$ is a real check on the metric, and it
costs one more fit.

\subsection{The reference-free acceptance}
\label{sec:identity}

There are no closed-form Solov'ev flux-surface averages. $\psi$ is elementary,
but $V'$ and the averages are integrals over a \emph{contour} of it, and those
contours have no elementary arc length. What a test fixture can supply is a
converged \emph{reference value}, and that distinction should be made wherever
the number is quoted.

So the sharpest check is one needing no reference at all. Taking the averaged
divergence identity with the vector $\gradbar\psi/R^{2}$ and using
$\dstar\psi = -F$,
\begin{equation}
	\frac{1}{V'} \frac{\mathrm{d}}{\mathrm{d}\psi}
		\left( V' \left\langle \frac{|\gradbar\psi|^{2}}{R^{2}} \right\rangle \right)
	= - \left\langle \frac{F}{R^{2}} \right\rangle .
	\label{eq:averaged-gs}
\end{equation}
Three averages checking each other with nothing but the equation.

Two things about \emph{measuring} with it, both of which caught a real problem.

\begin{quote}
	\textbf{The flux derivative must be Richardson-extrapolated.} A plain central
	difference carries its own second-order truncation, so the comparison is then
	floored by the \emph{instrument} rather than by the identity. Measured on a
	discrete field over a sixteenfold refinement, the plain column \emph{stops
	moving at all} while the extrapolated one keeps falling. A column that does
	not converge under mesh refinement is measuring the instrument, and an
	identity checked with that instrument would pass with a real defect
	underneath it.
\end{quote}

\begin{quote}
	\textbf{The step is a real choice and it is not monotone in either
	direction.} Too large and the extrapolation's own truncation dominates; too
	small and the difference of two nearly equal integrals divides the surfaces'
	own face-jump noise by the step and amplifies it. Measured, halving the step
	from near the optimum made the residual \emph{worse}. There is an optimum;
	find it on your own problem rather than assuming smaller is better.
\end{quote}

The right-hand side of \eqref{eq:averaged-gs} must be written with the $F$ the
solver is actually fed, and never re-expanded in terms of $p'$ and $gg'$: the
expanded form is source-specific and is re-derived by hand, so it is not
independent of the hand that derived it.

\subsection{Surfaces that cross the band}

Every node carries a flag saying whether it was reached by extension, and the
result carries the aggregate. \emph{Drop those surfaces before computing an error
norm or differencing two runs}: a surface with any band node is limited by the
extension's order rather than by the discretisation, and quoting one rate over
both populations hides precisely what the band flags exist to measure.

The contour builder marks a whole \emph{segment} rather than a node,
deliberately, because its nodes lie strictly between two accepted points and
cannot say for themselves. That over-reports by at most one segment at each end
of a band excursion and never under-reports --- and under-reporting is the
failure that matters, since it is band data presented as solved data.

\subsection{What the averages are not}

Not valid on an open surface: every rule here is periodic or closed. Not X-point
aware, for the reason \chref{ch:intro} gives. The safety factor \emph{carries no
evidence of its own} --- it is an algebraic combination of two quantities that
are measured, times a $g(\psi)$ the caller supplies, so it inherits their rates
and nothing more; what is worth pinning about it is its algebra and its constant,
which are exactly as easy to get wrong as anywhere and which no rate table would
see. And the arc length is \emph{not} an independent measurement of the metric,
since every weight already contains the arc-length element, so it recovers it
algebraically; what it checks is that the weights are built the way
\eqref{eq:fsa} says.

\section{The surfaces as a map from a disc}
\label{sec:disc}

The second deliverable is $R$ and $z$ as smooth functions of a flux label and an
in-surface angle: a map from a disc whose centre is the magnetic axis and whose
edge is the outermost surface, differentiable with respect to the label. That is
a fit over all the surfaces at once, and it is what a code needing the geometry as
a \emph{function} reads.

\subsection{The radial coordinate is the square root of the normalised flux}

\begin{quote}
	\textbf{Getting this wrong is silent.}
\end{quote}

$\psi$ has a \emph{quadratic} extremum at the magnetic axis --- a smooth function
with a non-degenerate critical point --- so the normalised flux behaves like the
square of the distance from the axis. The geometry is smooth in the distance,
therefore in the square root of the normalised flux, and it carries a square-root
branch point in the normalised flux itself.

Parametrise by the normalised flux directly and \emph{every} basis converges
algebraically against that branch point, worst near the axis, with nothing in a
convergence table to say why. The rate simply comes out below the design order
and looks like a discretisation problem. It is the same species of defect as a
wrong Jacobian: the answer converges, to the wrong rate, for a reason no
assertion is looking at.

The conversions between the two, and the chain factor between derivatives with
respect to each, are named functions rather than something a call site writes
out. The factor is one over twice the disc radius --- neither one nor twice the
radius, and both of those mistakes produce a quantity with the right units, the
right sign and the right qualitative shape.

That factor is unbounded at the axis, \emph{and that is the coordinate rather
than a defect}. Even for a representation perfectly smooth in Cartesian
coordinates, the derivative at fixed angle diverges at the axis whenever there is
any content in the first angular harmonic --- and those modes are the rigid shift
of a surface, so shifting the axis sideways at that rate as the surfaces shrink
onto a point is exactly what a Shafranov shift does. A consumer wanting something
finite there wants the derivative in the disc radius instead. Its growth as the
innermost surface approaches the axis is a conditioning number to read, not a
defect to fix.

\subsection{Why Zernike}

The Zernike polynomials on the unit disc pair a radial polynomial with an angular
harmonic, admissible exactly when the radial degree is at least the harmonic
order and the two differ by an even number. \emph{The index constraint is the
entire point of choosing the basis.}

A naive tensor product --- any polynomial in the radius times any Fourier mode in
the angle --- admits terms such as the square of the radius times the cosine of
the angle, which in Cartesian coordinates is $x\sqrt{x^{2}+y^{2}}$, whose second
derivative has no limit at the origin: it reads a different number down every
ray. Such a basis makes the centre of the disc a coordinate singularity that
every consumer has to special-case, and the special case is where the bugs live.

The constraint excludes exactly those terms, and what survives is the statement
worth having:

\begin{quote}
	\textbf{Every admissible Zernike mode is a bivariate polynomial in $(x,y)$.}
	So a Zernike expansion is a polynomial in Cartesian coordinates, the magnetic
	axis is an ordinary interior point of it, and derivatives there are whatever
	the polynomial says they are. No limit to take, no ray to choose, no branch
	--- at precisely the place where the surfaces shrink to a point, the weight
	in \eqref{eq:fsa} diverges, and the in-surface angle is undefined.
\end{quote}

The regularity condition this expresses is Lewis and Bellan's
\cite{LewisBellan1990}; the Fourier--Zernike basis satisfies it inherently, which
is the whole argument for it, and it is a better argument than that other codes
use it.

A consequence comes free and exactly: every mode with a non-zero harmonic carries
a positive power of the radius, so all of them vanish at the centre and what is
left there is one number with no angle in it. A Zernike fit therefore puts every
surface's collapse point at the same place \emph{by construction}, whatever the
coefficients are and however badly the fit was conditioned. What is not free is
whether that point is the magnetic axis: a sample set has a hole in the middle,
since nobody traces the surface at zero radius, so the fitted axis is an
extrapolation into that hole and is measured against the critical-point finder
rather than assumed.

\subsection{The angle is a label, and the obvious choice is the wrong one}

\begin{quote}
	\textbf{The geometric poloidal angle is inadmissible, and everything else in
	this section is downstream of that.}
\end{quote}

Labelling points on each surface by their geometric angle about the axis is the
obvious choice: it is what the angle parametrisation of
\S\ref{sec:angleparam} produces, it is
consistent across surfaces because the axis is one point, and it needs no extra
information. It is also, for any surface that is not a \emph{circle}, a
parametrisation in which the disc map is not smooth at the axis --- so no basis
that is smooth there can converge against it.

The argument is exact and takes one line. A function smooth at the origin has an
expansion whose coefficient of the first power of the radius carries exactly
\emph{one} angular harmonic. Take nested similar ellipses labelled by geometric
angle, and that coefficient carries the first, third, fifth and every subsequent
odd harmonic for ever --- every one of them beyond the first being a mode the
index constraint excludes, precisely because it is not smooth there. The
parametrisation puts content exactly where the basis refuses to look. Measured on
the ellipses themselves, where the answer is known exactly, the family is a
degree-one map under the ellipse's own parameter and is reproduced to round-off
at degree two, while the \emph{identical points} relabelled by geometric angle
are wrong at degree two and still wrong at degree twenty. Same points, same
basis, same solver; only the label differs.

\emph{The repair needs no field.} Near the axis every equilibrium's surfaces are
ellipses, since $\psi$ has a non-degenerate critical point there and the level
sets of a quadratic form are ellipses. Two numbers describe such an ellipse up to
scale --- how it is turned and how flat it is --- and both are recoverable from
the innermost traced surface alone. Relabelling by that ellipse turns geometric
angles into a label under which the leading term is smooth, and it is worth
orders of magnitude in the worst fit error at no cost in conditioning. It is an
\emph{exact reparametrisation} rather than an approximation: the map from
geometric angle to label is a fixed bijection of the circle, so a caller asking
where a surface is at a given geometric angle converts and gets the right point.
What is approximate is only the claim that the label makes the map smooth, and
that claim is exact at leading order and no better.

And that is a structural limit rather than a tuning problem. A relabelling that
does not depend on the radius has exactly one function's worth of freedom, so it
can fix exactly \emph{one order}: it makes the leading coefficient a single
harmonic and leaves everything above it carrying whatever the shaping puts there.
Pull the innermost traced surface in towards the axis and an algebraic tail
appears that no degree removes.

\subsection{The gauge-free fit}

Everything above pins the angle as an \emph{input}: each sample carries an angle,
so the fit is asked to put a particular point at a particular label. A code that
\emph{solves} for its coefficients is free to slide the angle to whatever its
basis can represent, and the DESC authors report exactly that of their own solver
--- it enforces no poloidal angle constraint and ends up finding a good
representation over the course of the optimisation \cite{PaniciDESC2023}. A
post-hoc fit cannot do that.

\emph{So ask for less.} Require each disc node only to \emph{land on the right
surface}: minimise, over the coefficients, the squared mismatch between the
normalised flux at the fitted position and the normalised flux the node belongs
on. Nothing says where \emph{along} its surface a node must sit, so the angle is
free and the truncated basis may choose it. The input to this is therefore a node
that \emph{carries no position} --- that is the entire difference from a sample.

It needs no force balance and no second physics solver, because {\meq} already
has $\psi$; and the Jacobian it needs is the gradient of the normalised flux,
which is the solved flux once more. Measured on nested ellipses, where the answer
is known exactly, it reaches round-off at low degree where the prescribed-angle
fit never converges at all --- and that is the theoretically right answer rather
than a lucky one, since for similar nested ellipses the family genuinely is a
degree-one map under the correct angle. On a discrete field it removes the
algebraic tail and the coefficient envelope resumes geometric decay.

Three things about it a consumer should know.

\emph{The residual is scaled into metres, and that is not cosmetic.} The
normalised flux is dimensionless and its gradient varies by orders across the
disc, so a least-squares problem in the raw flux residual weights outer surfaces
against inner ones by whatever that gradient happens to be. Dividing each row by
it makes the residual the normal distance from the node to its surface --- which
is the quantity a reader cares about, and is itself gauge invariant.

\emph{The warm start is not an optimisation.} The surface residual is minimised
by \emph{any} map onto the right surfaces, including folded ones and ones that
traverse a surface twice, so the basin matters. The prescribed-angle fit is what
puts the iterate in the right one; starting from zero coefficients is not
supported and would not be meaningful.

\emph{The freedom must be constrained, and the expectation about how was wrong.}
A tangential slide is invisible to the linearised system --- moving a node along
its own surface changes no residual --- so the Gauss--Newton matrix was expected
to be rank deficient by about half its columns. Measured, it is not, and the
answer depends on the field: on nested ellipses the exactly-null directions are a
handful, and on a realistic equilibrium there are \emph{none at all}. What both
have instead is a long soft tail with no gap in it. The reason is that a
tangential slide is exactly null only when the displaced map is still \emph{in}
the truncated space, and generally it is not --- so the gauge is not a subspace
to be projected out; it is a direction in which the problem is merely very soft,
and the floor deciding which directions to move has to be chosen rather than read
off a gap. And the soft tail is enough on its own: removing the gauge on a field
with no exact null space still leaves the map folded and the fit orders worse.
The default prefers the smallest step among those satisfying the surface
constraint, which is a spectral condensation preference \cite{HirshmanBreslau1998}
arrived at for free and with no weights to choose.

\begin{quote}
	\textbf{A fit with a beautiful residual and a folded map is the quiet wrong
	answer this machinery is most exposed to.} Nodes may bunch, cross and turn
	the map over while every residual stays perfect. The safeguard is the sign of
	the map's Jacobian, sampled over the disc --- and which sign is correct is not
	fixed and must not be assumed, since it depends on which way the angle runs.
\end{quote}

\subsection{A free angle changes what can be asserted}

Any check comparing the fit's position \emph{at a given angle} against a surface
traced at that angle is measuring precisely the freedom being granted. So a
gauge-free fit is accepted on \emph{gauge-invariant} properties instead: the
distance from each node to its surface, the fitted perimeter against the exact
one, and the sign of the map's Jacobian.

\begin{quote}
	\textbf{This costs the consumer nothing, and that is the point.} A coupled
	transport code reads flux-surface \emph{averages}, and an average does not
	know how its surface was parametrised --- \eqref{eq:fsa} is an integral over
	a curve. The deliverable was gauge invariant all along, which is what makes
	it legitimate to let the fit choose the angle.
\end{quote}

\subsection{What the representation is not}

Not a fit of open surfaces: every mode is periodic in the angle and regular at
the centre. Not an equilibrium solve --- nothing here enforces force balance,
$\psi$ is an input, and the only thing solved for is where the disc coordinates
put their points. It does not constrain the axis, because a node at zero radius
would ask for the flux to vanish, which is true there and would pin it for free
--- except that the gradient vanishes there too, so the row scaling divides by
zero and the linearisation carries no information about where to move. And
\emph{it does not know where its samples came from}, so it cannot flag a surface
as band-limited: mixing surfaces that crossed the band with surfaces that did not
fits two different accuracies into one expansion, which is a decision to be taken
deliberately rather than by default.

\section{Coupling: the geometry Jacobian}
\label{sec:jacobian}

A coupled transport code solves its own nonlinear system, in which {\meq}'s
geometry appears as coefficients. If it uses Newton --- and it should, for all
the reasons \chref{ch:intro} gives --- it needs the derivative of every geometry
slot with respect to every degree of freedom of $\psi$. This section is the
warning that goes with that requirement, and it is the most important thing in
this chapter.

\subsection{The term that is easy to forget}

A geometry slot is a flux-surface average, and a flux-surface average is an
integral over a surface \emph{whose location is itself determined by $\psi$}.
So
\begin{equation}
	\frac{\partial}{\partial \psi_m} \langle f \rangle
		= \underbrace{\left\langle \frac{\partial f}{\partial \psi_m}
			\right\rangle}_{\text{the integrand moves}}
		\; + \; \underbrace{\bigl(\text{terms from the surface moving}
			\bigr)}_{\text{easy to forget}} .
	\label{eq:jacobian}
\end{equation}

\begin{quote}
	\textbf{Both terms are needed.} Dropping the second gives a Jacobian that is
	wrong but \emph{plausible} --- right sign, right order of magnitude, right
	qualitative behaviour.
\end{quote}

\subsection{Why the failure is invisible}

Here is the part that makes this worth a section rather than a footnote.

\begin{quote}
	\textbf{An error in a forward Jacobian does not produce a wrong answer, only
	slow Newton convergence.} A consumer that never assembles its Jacobian ---
	which is to say, a matrix-free or Jacobian-free code, which most modern
	transport solvers are --- sees a missing block as a solver taking more
	iterations, and never as a bad number. Defects of exactly this kind have
	survived passing regression suites for months in the codes {\meq} is designed
	to couple to.
\end{quote}

The corollary is a rule rather than a suggestion: \emph{the geometry derivative
must be finite-difference tested directly against the geometry itself, as a unit
test.} Convergence of the coupled solve is not evidence that it is right. This
is the same discipline {\meq} applies to its own source terms, where the
derivative with respect to $\psi$ is differenced against the source in the unit
layer and the assembled Jacobian is differenced against the assembled residual in
the convergence layer --- the right instinct, applied to the right quantity.

An adjoint raises the stakes further. There a missing block does \emph{not} merely
slow anything down: it produces a silently wrong gradient beside a perfectly good
objective, and an optimiser will follow it.

\subsection{What is supplied, and what is not}

\begin{quote}
	\textbf{{\meq} does not currently supply the geometry Jacobian.} What it
	supplies is everything the derivative is taken \emph{of}, arranged so that
	taking it is one piece of work rather than eight.
\end{quote}

That arrangement is deliberate and is worth recognising when the derivative is
written. The averages are one facility rather than a list of named functions,
precisely so that the derivative is taken of the facility. The chain factor
between the two radial coordinates of \S\ref{sec:disc} is a named function rather
than something written out at each call site, because it is the factor that gets
dropped and because a factor lost in this derivative shows up as a coupled run
converging to the wrong equilibrium rather than as anything failing. And the disc
representation supplies both radial derivatives --- the one that diverges at the
axis and the one that does not --- so that whichever the consumer needs is the
one it gets.

Two routes to \eqref{eq:jacobian} exist and neither is implemented. Differencing
the whole extraction costs one surface extraction per degree of freedom of $\psi$,
which is expensive but is straightforward and is testable against nothing but
itself. A shape derivative --- differentiating the integral analytically, with the
surface velocity obtained from the level-set condition --- has a closed form for
both terms and is far cheaper, and it is the route to prefer if the cost of the
first proves unacceptable. Whichever is taken, the finite-difference test above
is what stands between it and a coupled run that is subtly and silently wrong.

\subsection{Calling convention}

One practical note about the shape of a coupling, learned the hard way on the
other side. A geometry hook is often \emph{pointwise}: called once per physics
node, per residual evaluation, handed the whole flux vector each time. For a
flux-surface average that means locating the surface through the point and
integrating on it, per node, per residual --- which is not affordable done
naively. A consumer taking this route must cache per $\psi$: extract the surface
geometry once when the flux changes, and serve individual points from that. If
the pointwise signature is the wrong shape, and it may well be, that is worth
raising rather than working around.
